\NativeTitle{45}{期中综合素养评价（二）}{}{90}{5}{一 & 二 & 三 & 四 & 五}{ &  &  &  & }
\NativeSlot{277}{44}{%
\ExamGroup{一、填空。{\normalfont\kaishu\fontsize{10.8bp}{14bp}\selectfont （第5题4分，其余每空1分，共23分）}}

% p045_s01

\ExamQuestion{1.}{5}%
给算式\(200+370-295\div5\)添上括号，使它的运算顺序变成先算减法，再算加法，最后算除法：\mbox{（\hspace{1.8em}）}。\par
}
\NativeSlot{352.0}{36}{%
% p045_s02

\ExamQuestion{2.}{5}%
装修公司要给一个面积为96平方米的会议厅铺防滑地砖，每平方米的地砖和施工费用合计48元，完成这个会议厅的地砖铺设总共需要花费\mbox{（\hspace{1.8em}）}元。\par
}
\NativeSlot{405.5}{36}{%
% p045_s03

\ExamQuestion{3.}{5}%
钟面上\mbox{（\hspace{1.8em}）}时整和\mbox{（\hspace{1.8em}）}时整，时针和分针互相垂直。\par
}
\NativeSlot{431.5}{36}{%
% p045_s04

\ExamQuestion{4.}{5}%
小新用加法运算律写了一个等式：\(\square2+4\square=\square8+5\square\)。如果这两个加数的和是100，那么这两个加数分别是\mbox{（\hspace{1.8em}）}和\mbox{（\hspace{1.8em}）}。\par
}
\NativeSlot{487.0}{49}{%
% p045_s05

\ExamQuestion{5.}{5}%
根据运算律填一填。\par\noindent （1）\(175+264+225+136=(\square\bigcirc\square)\bigcirc(\square\bigcirc\square)\)。\par

% p045_s06

\ExamQuestion{（2）}{5}%
\(\square\times99+\square\times\square=a\times(\square\bigcirc1)\)。\par
}
\NativeSlot{572}{36}{%
% p045_s07

\ExamQuestion{6.}{5}%
印刷厂工人每小时能印刷300份报纸，每天工作8小时，一周（5个工作日）能印刷\mbox{（\hspace{1.8em}）}份报纸。\par
}
\NativeSlot{625}{68}{%
% p045_s08

\ExamQuestion{7.}{5}%
在圆圈里填上“＞”“＜”或“＝”。\(38\times(5+6)\ \bigcirc\ 38\times5+6\)，\allowbreak{}\(420\div21\ \bigcirc\ 420\div7\div3\)；\allowbreak{}\(78\times199\ \bigcirc\ 78\times200-1\)，\allowbreak{}\(74\times10\times2\ \bigcirc\ 74\times12\)；\allowbreak{}\(65\times24\ \bigcirc\ 65\times4\times20\)，\allowbreak{}\(107\times5\ \bigcirc\ 105\times7\)。\par
}
\NativeSlot{743}{36}{%
% p045_s09

\ExamQuestion{8.}{5}%
斑马线中的每一组白线可以看作互相\mbox{（\hspace{1.8em}）}的线段。\par
}
\NativeSlot{769.0}{36}{%
% p045_s10

\ExamQuestion{9.}{5}%
一个等腰三角形的木板，顶角是底角的4倍，底角是\mbox{（\hspace{1.8em}）}°。\par
}
